\documentclass[10pt,a4paper]{amsart}
\usepackage[margin=19mm]{geometry}
\usepackage{amsmath,amssymb,mathtools}
\usepackage[scr=rsfs,cal=euler]{mathalfa}
\usepackage{xfrac}
\usepackage[unicode,hidelinks,bookmarks=false]{hyperref}
\newcommand{\ie}{i.e.\ }
\newcommand{\cf}{cf.\ }
\newcommand{\resp}{respectively\ }
\newcommand{\Subsection}{\subsection}
\providecommand{\nobreakdash}{\nobreak\hbox{-}}
\setlength{\emergencystretch}{2em}
\multlinegap0pt
\usepackage{bm}
\usepackage{bbm}
\usepackage{dsfont}
\usepackage{xspace}
\usepackage{cancel}
\usepackage{mathtools}
\usepackage{amsthm}
\makeatletter
\@ifundefined{lemma}{\newtheorem{lemma}{Lemma}}{}
\@ifundefined{proposition}{\newtheorem{proposition}[lemma]{Proposition}}{}
\makeatother
\title[On the average size of $1$-nearly independent vertex sets in]{On the average size of $1$-nearly independent vertex sets in graphs}
\author{Audace A. V. Dossou-Olory, Eric O. Andriantiana}
\date{Source: arXiv:2510.23670v1 (CC BY 4.0)}
\begin{document}
\maketitle

\noindent\emph{Source and licence.} On the average size of $1$-nearly independent vertex sets in graphs. Version arXiv:2510.23670v1 (CC BY 4.0), distributed under the Creative Commons Attribution 4.0 International licence, as recorded in the arXiv licence metadata for this exact version and shown on the abstract page https://arxiv.org/abs/2510.23670v1. Full rights notice: licenses/arxiv-2510.23670-CC-BY-4.0.txt.\par\smallskip

\noindent\textit{Editorial scope.} Counted unit: the Proposition whose source label is \texttt{prop:1sur3} in the article (source0.tex lines 1131--1141, source label \texttt{prop:1sur3}), delivered together with the article's own complete original proof of it (source0.tex lines 1143--1193, one \texttt{proof} environment of 51 lines). No mathematics is written, repaired or re-derived: statement and proof are the article's own text line for line. Required context retained uncounted: G1UnionG2 (lines 190--198); av1Exp (lines 291--299); lem:help (lines 1070--1084). Every definition, display and label of those lines is retained unchanged. Omitted from this package and not redistributed: 1--189, 199--290, 300--1069, 1085--1130, 1142--1142, 1194--1227. Nothing inside a retained excerpt is omitted or altered: every retained body is delivered exactly as the article's lines are written, so no omission receipt of any kind accompanies this package. Citations: the retained excerpts carry no citation command at all, so no bibliography entry of the article is transcribed and no reference is redistributed. Reference adaptation: none was needed and none was made; every reference inside the delivered unit resolves inside it, so no unresolved reference or citation is produced. Printed slips kept literal: no printed word, symbol, hypothesis or number of the counted statement or of its proof was altered. Layout and class: the article's own class is \texttt{amsart} with packages including a4wide, pgf, tikz, amssymb, enumerate, amsmath, makecell, relsize, lineno, caption, subcaption, float, hyperref, mathrsfs; the wrapper is the lane's pinned \texttt{amsart} editorial wrapper at 10pt on A4, which restates the theorem-like environments the retained text needs and adds the article's own macro definitions that the retained text needs (none), with extra wrapper packages (bm, bbm, dsfont, xspace, cancel, mathtools). The delivered numbering is the unit's own. Assets: no retained span contains a figure environment, a graphics-inclusion command, a PDF-page inclusion, a table or a TikZ picture; no image file is copied into this package, the package declares no asset dependency and no image file is redistributed with it. Licence: version arXiv:2510.23670v1 of this article is distributed under the Creative Commons Attribution 4.0 International licence, as recorded in the arXiv licence metadata for this exact version and shown on the abstract page https://arxiv.org/abs/2510.23670v1. A full rights notice is delivered at licenses/arxiv-2510.23670-CC-BY-4.0.txt. Characters: every retained excerpt is pure ASCII, so the delivered engine, XeLaTeX with its default Latin Modern text font, reports no missing character.
\begin{lemma}\label{G1UnionG2}
Let $G_1$ and $G_2$ be two vertex disjoint graphs. We have
\begin{align}
\label{Eq:Er1}
  I_1(G_1 \cup G_2;x)&=I_1(G_1;x)I_0(G_2;x)+I_1(G_2;x)I_0(G_1;x)\,\text{ and }\\
  \label{Eq:Er2}
 S_1(G_1 \cup G_2)&=S_1(G_1)\sigma_0(G_2)+\sigma_1(G_1)S_0(G_2) + S_1(G_2)\sigma_0(G_1)+\sigma_1(G_2)S_0(G_1)  \,.
\end{align}
\end{lemma}

\begin{lemma}\label{av1Exp}
For every graph $G$, it holds that
\begin{align*}
I_1(G;x)&=x^2 \sum_{uv \in E} I_0(G-(N(v) \cup N(u)); x)\,,\\
\sigma_1(G)&=\sum_{uv \in E} \sigma_0\big(G-(N(v) \cup N(u))\big)\,,\\
S_1(G)&=\sum_{uv \in E} \Big(2\sigma_0\big(G-(N(v) \cup N(u))\big) +S_0\big(G-(N(v) \cup N(u))\big) \Big)\\
&=\sum_{uv \in E} \sigma_0\big(G-(N(v) \cup N(u))\big)\Big(2 +av_0\big(G-(N(v) \cup N(u))\big) \Big)\,.
\end{align*}
\end{lemma}

\begin{lemma}\label{lem:help}
For every graph $G=(V,E)$ that is not edgeless and every $uv \in E$, we have 
\begin{align*}
\frac{1}{2^{l_{uv}-2}+2^{l_{uv}-m_u}+2^{l_{uv}-m_v}+1}\leq \frac{\sigma_0(G-(N[v]\cup N[u]))}{\sigma_0(G)}\leq 1- \frac{\sigma_0(G-v)}{\sigma_0(G)}\,,
\end{align*}
where
\begin{align*}
l_{uv}=|N(v)\cup N(u)|\,, \quad m_u=|N[u]|\,, \quad m_v=| N[v]|\,.
\end{align*}
In particular, 
\begin{align*}
\frac{\sigma_0(G-(N[v]\cup N[u])}{\sigma_0(G)}\geq \frac{1}{(3\cdot 2^{l_{uv}-2}+1)}\,.
\end{align*}

\end{lemma}

\begin{proposition}\label{prop:1sur3}
Let $G_1$ and $G_2$ be two vertex disjoint graphs. We have
\begin{align*}
\frac{ \sigma_1(G_1 \cup G_2)}{ \sigma_0(G_1 \cup G_2)}=\frac{ \sigma_1(G_1)}{ \sigma_0(G_1)}+\frac{ \sigma_1(G_2)}{ \sigma_0(G_2)}\,.
\end{align*}
Moreover, for every graph $G=(V,E)$ with no isolated vertex and maximum degree $2$,
\begin{align*}
\frac{\sigma_1(G)}{\sigma_0(G)}\geq \frac{1}{3}\,.
\end{align*}
Equality holds for $P_2$.
\end{proposition}

\begin{proof}
The first identity immediately follows from Lemma~\ref{G1UnionG2}:
\begin{align*}
\frac{ \sigma_1(G_1 \cup G_2)}{ \sigma_0(G_1 \cup G_2)}=\frac{\sigma_1(G_1) \sigma_0(G_2) + \sigma_1(G_2) \sigma_0(G_1)}{\sigma_0(G_1) \sigma_0(G_2)}\,.
\end{align*}
Now let $G=(V,E)$ be a graph with no isolated vertex and maximum degre $2$. All the connected components of $G$ are paths and/or cycles. If $P_2$ is such a component, then set $G_1:=P_2$ and $G_2:=G\backslash V(G_1)$ to obtain
\begin{align*}
\frac{ \sigma_1(G)}{ \sigma_0(G)}=\frac{ \sigma_1(P_2)}{ \sigma_0(P_2)}+\frac{ \sigma_1(G_2)}{ \sigma_0(G_2)}\geq \frac{ \sigma_1(P_2)}{ \sigma_0(P_2)}=\frac{1}{3}\,.
\end{align*}
Otherwise, 
\begin{align*}
l_{uv}:=|N(v)\cup N(u)|\in \{3,4\}
\end{align*}
for all $uv\in E$. Here we distinguish two cases:
\begin{itemize}
\item Case~1: $|E|\geq 5$.\\
We can partition $E$ as $E=E_4 \cup E_3$, where
\begin{align*}
E_4=\{uv\in E:~ l_{uv}=4 \}\quad \text{and} \quad E_3=\{uv\in E:~ l_{uv}=3 \}\,.
\end{align*}
We combine (see Lemma~\ref{lem:help})
\begin{align*}
\frac{\sigma_0(G-(N[v]\cup N[u])}{\sigma_0(G)}\geq \frac{1}{(3\cdot 2^{l_{uv}-2}+1)}\,.
\end{align*}
and (see~Lemma~\ref{av1Exp})
\begin{align*}
\frac{\sigma_1(G)}{\sigma_0(G)}=\sum_{uv \in E} \frac{\sigma_0\big(G-(N(v) \cup N(u))\big)}{\sigma_0(G)}
\end{align*}
to obtain
\begin{align*}
\frac{\sigma_1(G)}{\sigma_0(G)}&\geq \sum_{uv \in E_4} \frac{1}{(3\cdot 2^{4-2}+1)} ~ + ~\sum_{uv \in E_3} \frac{1}{(3\cdot 2^{3-2}+1)}=\frac{|E_4|}{13}+\frac{|E_3|}{7}\\
& \geq \frac{5-|E_3|}{13}+\frac{|E_3|}{7}=\frac{35+6|E_3|}{91}> \frac{1}{3}\,.
\end{align*}
\item Case~2: $1\leq |E|\leq 4$.\\
Here all the connected components of $G$ must belong to the set
\begin{align*}
Q:=\{P_5,C_4,P_4,C_3,P_3\}\,,
\end{align*}
while easy calculations yield
\begin{align*}
&\frac{\sigma_1(P_5)}{\sigma_0(P_5)}=\frac{10}{13}\,, \quad \frac{\sigma_1(C_4)}{\sigma_0(C_4)}=\frac{4}{7}\,, \quad \frac{\sigma_1(P_4)}{\sigma_0(P_4)}=\frac{5}{8}\,, \\
&\frac{\sigma_1(C_3)}{\sigma_0(C_3)}=\frac{3}{4}\,, \quad \frac{\sigma_1(P_3)}{\sigma_0(P_3)}=\frac{2}{5}\,.
\end{align*}
All these values are greater than $1/3$. Now $G=G_1 \cup G_2$ for some $G_1\in Q$. We have
\begin{align*}
\frac{ \sigma_1(G_1 \cup G_2)}{ \sigma_0(G_1 \cup G_2)}=\frac{ \sigma_1(G_1)}{ \sigma_0(G_1)}+\frac{ \sigma_1(G_2)}{\sigma_0(G_2)} > \frac{1}{3}\,,
\end{align*}
which concludes the proof.
\end{itemize}

\end{proof}

\end{document}
